\documentclass[11pt]{article}

\usepackage[margin=1in]{geometry}
\usepackage{times}
\usepackage{microtype}
\usepackage{enumitem}
\usepackage{hyperref}

\title{\vspace{-0.8em}
CSE 5349 Project Proposal:\\
Finding and Understanding \texttt{unsafe}-Related Failures in Real Rust Crates\\
with Hotspot Mining, Miri, and Fuzzing
\vspace{-0.6em}}
\author{TBD}
\date{}

\begin{document}
\maketitle
\vspace{-1.0em}

\section*{Summary}
We will run a case study on several real-world, input-facing Rust crates (e.g., a parser, decoder, or deserializer) to see how their \texttt{unsafe} code can fail. Specifically, we will:
\begin{enumerate}[leftmargin=*]
  \item Map \texttt{unsafe} ``hotspots'' in each crate and its dependency graph.
  \item Run each crate's test suite under \textbf{Miri} to catch undefined behavior (UB) from violated \texttt{unsafe} assumptions.
  \item Write 1--3 \textbf{coverage-guided fuzz targets} per crate around input-facing APIs to hunt for panics, crashes, and edge cases.
  \item Write up where \texttt{unsafe} appears, what broke (if anything), and how Miri and fuzzing each performed across the crates.
\end{enumerate}

\section*{Motivation}
Rust prevents many memory safety bugs at compile time, but \texttt{unsafe} is a deliberate escape hatch. Code inside \texttt{unsafe} blocks can do fast, low-level work or call into C, at the cost of extra invariants the compiler will not check. Evans et al.\ found that most crates avoid writing \texttt{unsafe} directly but still pull it in through dependencies~\cite{EvansICSE20}. Astrauskas et al.\ studied how programmers actually use \texttt{unsafe} and how often those uses are sound~\cite{AstrauskasOOPSLA20}. Qin et al.\ cataloged real memory and thread safety bugs in Rust programs, many tied to \texttt{unsafe} regions~\cite{QinPLDI20}. On the tooling side, RUDRA~\cite{BaeSOSP21} finds \texttt{unsafe}-related bugs at ecosystem scale, and Miri~\cite{JungPOPL26} interprets Rust at the MIR level to detect UB on executed paths.

We want to bring these ideas down to a handful of real crates and ask a concrete question: for each library, where does \texttt{unsafe} risk actually sit, and what do we learn by throwing Miri and a fuzzer at it in a controlled setup?

\section*{Problem statement}
Given a set of real Rust crates and their dependency closures:
\begin{itemize}[leftmargin=*]
  \item It is hard to tell where the most security-relevant \texttt{unsafe} code lives---whether it is in a crate itself or buried in a dependency.
  \item Miri and fuzzing catch different kinds of problems (UB on executed paths vs.\ panics from unexpected inputs), but for any given crate it is not obvious which tool will find what, and comparing across multiple crates gives us a broader picture.
  \item We want a project that is small enough for a semester but still goes beyond building an application.
\end{itemize}

\section*{Goals}
\textbf{G1:} Produce an \texttt{unsafe} hotspot map for each target crate, including dependencies.\\
\textbf{G2:} Run Miri and fuzzing on each crate. Record any findings with minimal reproducers.\\
\textbf{G3:} Compare what each tool found and what it missed, across the target crates.

\section*{Approach}
\subsection*{Target selection}
Choose several crates, each satisfying:
\begin{itemize}[leftmargin=*]
  \item has an input-driven API suitable for fuzzing (e.g., parse/decode/from\_bytes),
  \item has a test suite that runs reliably on a pinned toolchain,
  \item contains some \texttt{unsafe}, either directly or through a key dependency.
\end{itemize}
We will shortlist candidates and finalize based on a quick \texttt{unsafe} scan.

\subsection*{Step A: Hotspot mining}
Run \texttt{cargo-geiger} on each target crate to count and locate \texttt{unsafe} usage across its dependency graph. Output: a ranked list of the modules and dependencies with the most \texttt{unsafe}, per crate.

\subsection*{Step B: Miri on existing tests}
Run each crate's tests under Miri. If Miri reports UB, minimize it into a short reproducer and trace it back to the responsible code. Even a single Miri finding shows that an \texttt{unsafe} assumption can actually be violated under normal test inputs.

\subsection*{Step C: Coverage-guided fuzzing}
Write 1--3 fuzz targets per crate with \texttt{cargo-fuzz}, pointed at the most input-facing entry points (e.g., parse/decode). For each target:
\begin{itemize}[leftmargin=*]
  \item seed the corpus with a handful of valid and near-valid inputs,
  \item run fuzzing for a fixed time budget,
  \item triage crashes and panics, minimize inputs,
  \item record stack traces and likely root cause (whether \texttt{unsafe}-related or not).
\end{itemize}

\subsection*{Step D: Reporting}
Write a short case study report and a reproducible artifact (scripts + instructions to rerun everything).

\section*{Implementation plan}
\begin{itemize}[leftmargin=*]
  \item A runner script that reproduces:
  \begin{itemize}[leftmargin=1.2em]
    \item the \texttt{cargo-geiger} scan and hotspot summary,
    \item \texttt{cargo miri test} for selected test targets,
    \item \texttt{cargo fuzz run} for each fuzz target with time budget.
  \end{itemize}
  \item 1--3 fuzz harnesses.
\end{itemize}

\section*{Evaluation}
This is a case study, so we evaluate on clarity and evidence:
\begin{enumerate}[leftmargin=*]
  \item Hotspot summary: where \texttt{unsafe} lives, ranked by concentration.
  \item Tool outcomes: did Miri find UB? Did fuzzing find crashes or panics within the time budget?
  \item For each issue found (if any): a minimal reproducer and a short root-cause explanation.
  \item A comparison of Miri vs.\ fuzzing across the target crates, with practical notes on effort and limitations.
  \item Reproducibility: scripts and instructions for rerunning the whole study.
\end{enumerate}
If we find no bugs, the project is still useful as long as it produces clear hotspot maps and an honest Miri-vs-fuzz comparison across real crates.

\section*{Questions for Mike}
\begin{itemize}[leftmargin=*]
  \item Do you have target crate suggestions? (Mid-sized, input-facing, some \texttt{unsafe}, decent tests.)
  \item Is a case study with a tool comparison a sufficient outcome, even if we find nothing exploitable?
\end{itemize}

\vspace{0.5em}
\noindent\textbf{References:} empirical unsafe usage and issues \cite{EvansICSE20,AstrauskasOOPSLA20,QinPLDI20}; ecosystem bug finding \cite{BaeSOSP21}; semantics-aware UB detection \cite{JungPOPL26}.

\begin{thebibliography}{9}

\bibitem{EvansICSE20}
Ana Nora Evans, Bradford Campbell, and Mary Lou Soffa.
\newblock \emph{Is Rust Used Safely by Software Developers?}
\newblock In \emph{ICSE}, 2020.
\newblock \url{https://www.cs.virginia.edu/~bjc8c/papers/evans20rust.pdf}

\bibitem{AstrauskasOOPSLA20}
Vytautas Astrauskas, Darrell Matheja, Federico Poli, and Alexander J. Summers.
\newblock \emph{How Do Programmers Use Unsafe Rust?}
\newblock \emph{Proc.\ ACM Program.\ Lang. (OOPSLA)}, 2020.
\newblock \url{https://pm.inf.ethz.ch/publications/AstrauskasMathejaMuellerPoliSummers20.pdf}

\bibitem{QinPLDI20}
Boqin Qin, Yilun Chen, Zeming Yu, Linhai Song, and Yiying Zhang.
\newblock \emph{Understanding Memory and Thread Safety Practices and Issues in Real-World Rust Programs.}
\newblock In \emph{PLDI}, 2020.
\newblock \url{https://cseweb.ucsd.edu/~yiying/RustStudy-PLDI20.pdf}

\bibitem{BaeSOSP21}
Yechan Bae, Youngsok Kim, Ammar Askar, Jungwon Lim, and Taesoo Kim.
\newblock \emph{RUDRA: Finding Memory Safety Bugs in Rust at the Ecosystem Scale.}
\newblock In \emph{SOSP}, 2021.
\newblock \url{https://gts3.org/assets/papers/2021/bae%3Arudra.pdf}

\bibitem{JungPOPL26}
Ralf Jung, Benjamin Kimock, Christian Poveda, Eduardo Sanchez Munoz, Oli Scherer, and Qian Wang.
\newblock \emph{Miri: Practical Undefined Behavior Detection for Rust.}
\newblock \emph{Proc.\ ACM Program.\ Lang. (POPL)}, 2026.
\newblock \url{https://research.ralfj.de/papers/2026-popl-miri.pdf}

\end{thebibliography}

\end{document}